\documentclass[11pt,a4paper]{article}

\usepackage[utf8]{inputenc}
\usepackage[T1]{fontenc}
\usepackage{amsmath,amssymb}
\usepackage{physics}
\usepackage{braket}
\usepackage{hyperref}
\usepackage[margin=1in]{geometry}
\usepackage{enumitem}
\usepackage{xcolor}
\usepackage{listings}

\lstset{
  basicstyle=\ttfamily\footnotesize,
  breaklines=true,
  showstringspaces=false,
  keywordstyle=\color{blue},
  commentstyle=\color{gray},
}

\newcommand{\Dslash}{D}
\newcommand{\Slambda}{S_\lambda}
\newcommand{\Szero}{S}
\newcommand{\Sresp}{S_{\mathrm{resp}}}
\newcommand{\Ocur}{\mathcal{O}}
\newcommand{\gcur}{\Gamma_{\mathrm{cur}}}
\newcommand{\gsrc}{\Gamma_{\mathrm{src}}}
\newcommand{\gsink}{\Gamma_{\mathrm{sink}}}
\newcommand{\gfive}{\gamma_5}
\newcommand{\Phiq}{\Phi_{\mathbf q}}
\newcommand{\Ctwo}{C_2}
\newcommand{\Cthree}{C_3}
\newcommand{\Cresp}{C_{\mathrm{resp}}}
\newcommand{\Cresptwo}{C_{\mathrm{resp}}^{(2)}}
\newcommand{\Rsum}{R_{\mathrm{sum}}}
\newcommand{\pyquda}{\texttt{PyQUDA}}

\title{\textbf{Pion Current Background-Response Diagnostics}\\[4pt]
\large Linearized Current Response and Current-Current Response Propagators}
\author{Pyquda\_Measurement Collaboration}
\date{\today}

\begin{document}
\maketitle
\tableofcontents

%=============================================================================
\section{Purpose}
%=============================================================================

This note documents the pion current background-response diagnostics implemented
in
\begin{center}
\texttt{pyquda\_measurement\_utils/pion\_current\_background\_response\_vibe\_develop.py}
\end{center}
and
\begin{center}
\texttt{application/EMFF\_pion\_background\_response/perlmutter/}.
\end{center}
The production utility contains only inserted-source construction, response
inversions, the response pion-C2 contraction, and relative-to-absolute
projector conversion. Tau-window selection, source-relative rolling,
ratio/channel extraction, explicit-EMFF summation, and HDF5 writers are
centralized in
\texttt{application/analysis\_helper/pion\_current\_response\_analysis.py}.

The goal is to study background-field/Feynman-Hellmann routes to pion current
matrix elements while \emph{not} modifying the QUDA Dirac operator and while
\emph{not} performing inversions with a finite external field
\(\Dslash+\lambda\Ocur\).  Instead, the current implementation constructs
linearized response propagators using ordinary \(\Dslash^{-1}\) inversions.

This is not yet the full finite-\(\lambda\) background-field method.  It is the
linearized method at \(\lambda=0\), which is mathematically equivalent to a
summed current insertion.

The pion electromagnetic form-factor (EMFF) response is the first concrete
application, using the temporal vector current \(V_4\).  The same source module
also contains a minimal second-order current-current response building block for
future hadronic-tensor or Compton-amplitude style diagnostics.

%=============================================================================
\section{Two Related But Distinct Methods}
%=============================================================================

\subsection{Finite-\(\lambda\) background-field method}

The most literal background-field method modifies the action or Dirac operator,
\begin{equation}
  \Dslash_\lambda = \Dslash + \lambda \Ocur,
\end{equation}
then computes ordinary two-point functions in this modified theory,
\begin{equation}
  C_\lambda(t)
  =
  \braket{
    O_\pi(t) O_\pi^\dagger(0)
  }_\lambda .
\end{equation}
At large Euclidean time,
\begin{equation}
  C_\lambda(t)
  =
  Z_\lambda e^{-E_\lambda t}
  + \cdots .
\end{equation}
The desired matrix element is obtained from the energy response,
\begin{equation}
  \left.\frac{\partial E_\lambda}{\partial\lambda}\right|_{\lambda=0}.
\end{equation}
This is the usual ``extract the energy shift'' picture of the background-field
method.

\subsection{Linearized response at \(\lambda=0\)}

The method currently implemented in this repository does not invert
\(\Dslash+\lambda\Ocur\).  Instead, it expands the propagator,
\begin{equation}
  \Slambda
  =
  (\Dslash + \lambda \Ocur)^{-1}
  =
  \Szero - \lambda \Szero \Ocur \Szero + O(\lambda^2).
\end{equation}
Thus,
\begin{equation}
  \left.
  \frac{\partial \Slambda}{\partial\lambda}
  \right|_{\lambda=0}
  =
  - \Szero \Ocur \Szero.
\label{eq:dSdlambda}
\end{equation}

The code constructs the object
\begin{equation}
  \Sresp
  =
  \Szero \Ocur \Szero
  =
  \Dslash^{-1} \Ocur \Szero,
\label{eq:Sresponse_code_sign}
\end{equation}
using ordinary inversions.  Therefore, relative to the derivative in
Eq.~\eqref{eq:dSdlambda}, the physical finite-difference sign is
\begin{equation}
  \left.
  \frac{\partial \Slambda}{\partial\lambda}
  \right|_{\lambda=0}
  =
  -\Sresp.
\end{equation}
The positive-sign \(\Sresp\) convention is useful because it can be compared
directly against the explicit summed three-point function.

\subsection{Why the linear response looks like a summed three-point function}

Expanding the two-point function in the external source gives
\begin{equation}
  \left.
  \frac{\partial C_\lambda(t)}{\partial\lambda}
  \right|_{\lambda=0}
  =
  -\sum_\tau
  \braket{
    O_\pi(t)\,
    \Ocur(\tau)\,
    O_\pi^\dagger(0)
  },
\label{eq:dC_summed_3pt_sign}
\end{equation}
up to the sign convention chosen in the action perturbation.  In the code-sign
convention of Eq.~\eqref{eq:Sresponse_code_sign}, the response contraction is
compared to
\begin{equation}
  \sum_\tau \Cthree(t,\tau).
\end{equation}
This explains why a linearized background-response calculation has analysis
features similar to a summed three-point function: it is the derivative
\(\partial C_\lambda/\partial\lambda|_0\), not a finite-\(\lambda\) energy fit.

%=============================================================================
\section{Pion Current Kinematics}
%=============================================================================

For the pion EMFF application, the current matrix element is defined by
\begin{equation}
  \braket{\pi(p_f)|J_\mu(0)|\pi(p_i)}
  =
  (p_i+p_f)_\mu F_\pi(Q^2),
  \qquad
  q=p_f-p_i .
\end{equation}
The lattice implementation uses
\begin{equation}
  \mathbf p_i = \mathbf p_f - \mathbf q_{\mathrm{ext}}.
\end{equation}
The explicit three-point correlator is
\begin{align}
  \Cthree^\Gamma(\mathbf q,\tau;\mathbf p_f,t_{\mathrm{sep}})
  &=
  \sum_{\mathbf x}
  e^{-i\mathbf q\cdot(\mathbf x-\mathbf x_0)}
  \tr\left[
    S_{\mathrm{seq,anti}}(\mathbf x,\tau;\mathbf p_f,t_{\mathrm{sep}})
    \Gamma
    S_q(\mathbf x,\tau;\mathbf x_0,0)
    \gsrc
  \right],
\label{eq:explicit_emff_c3}
\end{align}
where the fixed-sink sequential propagator has already absorbed the sink
momentum \(\mathbf p_f\), the sink time \(t_{\mathrm{sep}}\), and the sink
interpolator \(\gsink\).  The standard EMFF choice is the temporal vector
current,
\begin{equation}
  \Gamma = \gamma_T,
\end{equation}
labeled \(\texttt{T}\) in the code.

%=============================================================================
\section{Current-Inserted Source}
%=============================================================================

For a local current with momentum transfer \(\mathbf q\),
\begin{equation}
  \Ocur_{\Gamma,\mathbf q}(x)
  =
  \Gamma\,\Phi_{\mathbf q}(x-\mathbf x_0),
\end{equation}
where \(\Phi_{\mathbf q}\) is produced by \texttt{MomentumPhase} using the same
phase convention as the ordinary pion current three-point code.

Given an ordinary source-to-current propagator
\begin{equation}
  S(x,x_0),
\end{equation}
the all-time current-inserted source is
\begin{equation}
  \eta_{\Gamma,\mathbf q}(x)
  =
  \Gamma\,\Phi_{\mathbf q}(x-\mathbf x_0)\,S(x,x_0).
\label{eq:inserted_source_alltime}
\end{equation}
The response propagator is then obtained by a standard inversion,
\begin{equation}
  \Sresp(y,x_0;\Gamma,\mathbf q)
  =
  \sum_x S(y,x)\,
  \Gamma\,\Phi_{\mathbf q}(x-\mathbf x_0)\,
  S(x,x_0).
\label{eq:response_prop}
\end{equation}

In code, Eq.~\eqref{eq:inserted_source_alltime} is built by
\begin{center}
\texttt{build\_local\_current\_inserted\_source(...)}
\end{center}
and Eq.~\eqref{eq:response_prop} by
\begin{center}
\texttt{invert\_local\_current\_response\_propagator(...)}.
\end{center}

%=============================================================================
\section{Response Correlator}
%=============================================================================

The response propagator can be contracted like the forward quark line in a pion
two-point function.  With an ordinary antiquark line \(S_{\mathrm{anti}}\),
\begin{align}
  \Cresp(t;\mathbf p_f,\mathbf q)
  &=
  \sum_{\mathbf y}
  e^{-i\mathbf p_f\cdot(\mathbf y-\mathbf x_0)}
  \tr\left[
    S_{\mathrm{anti}}(\mathbf y,t;\mathbf x_0,0)\,
    \gsink\,
    \Sresp(\mathbf y,t;\mathbf x_0,0;\Gamma,\mathbf q)\,
    \gsrc
  \right].
\label{eq:Cresponse}
\end{align}

For the code-sign convention,
\begin{equation}
  \Cresp(t_{\mathrm{sep}};\mathbf p_f,\mathbf q)
  =
  \sum_\tau
  \Cthree^\Gamma(\mathbf q,\tau;\mathbf p_f,t_{\mathrm{sep}}),
\label{eq:Cresponse_equals_summed_C3}
\end{equation}
provided the same current gamma, source/sink interpolators, momentum phases, and
smearing conventions are used.

The first validated case is
\begin{equation}
  \mathbf p_f=\mathbf 0,\qquad
  \mathbf q=\mathbf 0,\qquad
  \mathbf p_i=\mathbf 0,\qquad
  \Gamma=\gamma_T,\qquad
  \gsrc=\gsink=\gamma_5.
\end{equation}
On the S8T32 test gauge, the diagnostic found
\begin{align}
  \sum_\tau C_3(\tau)
    &= 327.27206385480486 + 0.1342484007638715\,i,\\
  C_{\mathrm{resp}}
    &= 327.27206385480486 + 0.13424840076386432\,i,
\end{align}
with relative difference
\begin{equation}
  2.2\times 10^{-17}.
\end{equation}

%=============================================================================
\section{Energy-Shift Interpretation}
%=============================================================================

Although Eq.~\eqref{eq:Cresponse_equals_summed_C3} looks like a summed
three-point function, it is still directly related to the background-field
energy shift.  Suppose
\begin{equation}
  C_\lambda(t)
  =
  Z_\lambda e^{-E_\lambda t}.
\end{equation}
Expanding around \(\lambda=0\),
\begin{align}
  Z_\lambda &= Z_0 + \lambda Z'_0 + O(\lambda^2),\\
  E_\lambda &= E_0 + \lambda E'_0 + O(\lambda^2),
\end{align}
gives
\begin{equation}
  C_\lambda(t)
  =
  C_0(t)
  \left[
    1+\lambda
    \left(
      \frac{Z'_0}{Z_0} - t E'_0
    \right)
  \right]
  + O(\lambda^2).
\end{equation}
Therefore
\begin{equation}
  \frac{1}{C_0(t)}
  \left.
  \frac{\partial C_\lambda(t)}{\partial\lambda}
  \right|_{\lambda=0}
  =
  \frac{Z'_0}{Z_0} - t E'_0.
\label{eq:response_ratio_slope}
\end{equation}

Thus, in the linear-response method, the energy-shift derivative is extracted
from the slope of a response ratio, not from separate finite-\(\lambda\) energy
fits.  This is the same Feynman-Hellmann information, but represented at
\(\lambda=0\).

%=============================================================================
\section{Summed-Ratio Object}
%=============================================================================

Define
\begin{equation}
  \Rsum(t_{\mathrm{sep}})
  =
  \frac{
    \sum_\tau \Cthree(t_{\mathrm{sep}},\tau)
  }{
    \Ctwo(t_{\mathrm{sep}})
  }
  =
  \frac{
    \Cresp(t_{\mathrm{sep}})
  }{
    \Ctwo(t_{\mathrm{sep}})
  }.
\end{equation}
At large \(t_{\mathrm{sep}}\),
\begin{equation}
  \Rsum(t_{\mathrm{sep}})
  =
  A - t_{\mathrm{sep}} E'_0
  + \text{excited-state corrections},
\label{eq:summed_ratio_slope}
\end{equation}
where \(A\) absorbs overlap-factor response and contact-term effects.

For the finite-\(\lambda\) sign convention
\(\Dslash_\lambda=\Dslash+\lambda\Ocur\), the derivative of the propagator has
the sign in Eq.~\eqref{eq:dSdlambda}.  The diagnostic code stores both
\begin{center}
\texttt{response\_sign = +1}
\end{center}
for direct comparison to the explicit summed three-point function, and
\begin{center}
\texttt{finite\_difference\_derivative\_sign = -1}
\end{center}
for the corresponding finite-field derivative convention.

%=============================================================================
\section{Restricted Insertion-Time Windows}
%=============================================================================

The full background-field energy-shift picture corresponds to an all-time
insertion.  However, in the linearized response implementation one may restrict
the source-relative insertion time \(\tau_{\rm rel}\).  For a source at
absolute time \(t_0\), the projector is constructed only after applying
\begin{equation}
  \tau_{\rm abs}=(t_0+\tau_{\rm rel})\bmod N_t .
\end{equation}
Thus one writes
\begin{equation}
  \Sresp^{W_\tau}
  =
  \sum_{\tau\in W_\tau}
  \Dslash^{-1}
  \left[
    \delta_{t,\tau}\,
    \Gamma\Phi_{\mathbf q}(x-\mathbf x_0)S(x,x_0)
  \right].
\end{equation}
This produces
\begin{equation}
  \Cresp^{W_\tau}(t_{\mathrm{sep}})
  =
  \sum_{\tau\in W_\tau}
  \Cthree(t_{\mathrm{sep}},\tau).
\end{equation}

Useful choices include
\begin{align}
  W_\tau &= \text{all time slices},\\
  W_\tau &= \{0,\ldots,t_{\mathrm{sep}}\},\\
  W_\tau &= \{1,\ldots,t_{\mathrm{sep}}-1\},\\
  W_\tau &= \{\tau_{\min},\ldots,t_{\mathrm{sep}}-\tau_{\min}\}.
\end{align}
The restricted windows are not the pure finite-field energy-shift method.
They are analysis tools for the linearized method, useful for:
\begin{itemize}
  \item removing source/sink contact terms,
  \item comparing with explicit three-point functions in the same time window,
  \item studying excited-state contamination,
  \item checking the stability of summed-ratio slopes.
\end{itemize}

%=============================================================================
\section{Current Implementation}
%=============================================================================

\subsection{Source helper}

The helper
\begin{center}
\texttt{build\_local\_current\_inserted\_source(..., source\_time, tau\_relative\_list)}
\end{center}
constructs
\begin{equation}
  \eta(x)
  =
  \Gamma_{\mathrm{cur}}\Phi_{\mathbf q}(x-\mathbf x_0)S(x,x_0).
\end{equation}
If \texttt{tau\_relative\_list} is provided, each entry is mapped periodically
to absolute lattice time immediately before calling the time-projector helper.
If it is omitted, all time slices are included.

\subsection{Response inversion}

The helper
\begin{center}
\texttt{invert\_local\_current\_response\_propagator(...)}
\end{center}
performs the ordinary \(\Dslash^{-1}\) inversion of the current-inserted
source.  If a list \texttt{tau\_relative\_list} is given, the code first builds
the summed time-window source and then performs one inversion.  This deliberately
avoids storing per-\(\tau\) response propagators.

\subsection{Response contraction}

The helper
\begin{center}
\texttt{contract\_response\_pion\_2pt(...)}
\end{center}
contracts the response propagator with the ordinary antiquark line using the
same pion two-point contraction infrastructure as the rest of the pion code.

\subsection{Current-current response}

The same helper module also contains a minimal nested current-current response
prototype,
\begin{center}
\texttt{invert\_current\_current\_response\_propagator(...)}.
\end{center}
For two local currents \(O_1\) and \(O_2\), it constructs
\begin{equation}
  \Sresp^{(2)}
  =
  \Dslash^{-1} O_2 \Dslash^{-1} O_1 S.
\end{equation}
With
\begin{equation}
  O_i(x)
  =
  \Gamma_i \Phi_{\mathbf q_i}(x-\mathbf x_0),
\end{equation}
this corresponds to the ordered two-current insertion
\begin{equation}
  S^{(2)}_{21}
  =
  S O_2 S O_1 S,
\end{equation}
up to the sign convention chosen for finite-\(\lambda\) derivatives.  The
associated pion two-point-like response contraction is
\begin{equation}
  \Cresptwo(t_{\mathrm{sep}})
  =
  \sum_{\tau_2\in W_2}
  \sum_{\tau_1\in W_1}
  C_{JJ}(t_{\mathrm{sep}},\tau_2,\tau_1),
\end{equation}
where \(C_{JJ}\) is the pion correlator with two local current insertions on
the active quark line.  The total momentum injected by the two currents is
\begin{equation}
  \mathbf q_{\mathrm{tot}} = \mathbf q_1 + \mathbf q_2,
  \qquad
  \mathbf p_i = \mathbf p_f - \mathbf q_{\mathrm{tot}}.
\end{equation}
This is the direct local-current analogue of a two-current insertion into the
quark line.  It is intended as a diagnostic building block for current-current
and hadronic-tensor style measurements.  As with the first-order response, tau
windows are summed at the inserted-source level before inversion; the code does
not save per-\(\tau\) response-propagator caches.

\subsection{Diagnostic application}

The current first-order Perlmutter diagnostic application is
\begin{center}
\texttt{application/EMFF\_pion\_background\_response/perlmutter/}.
\end{center}
The default run uses the S8T32 test gauge and compares
\begin{equation}
  \Cresp(t_{\mathrm{sep}})
  \quad\text{against}\quad
  \sum_\tau \Cthree(t_{\mathrm{sep}},\tau)
\end{equation}
for \(t_{\mathrm{sep}}=2\), \(\Gamma=\gamma_T\), and zero external momentum.
The upgraded diagnostic can also run several \(t_{\mathrm{sep}}\), several
momentum transfers, and several current gamma matrices in a single job.  The
HDF5 output stores one record for each
\((\Gamma,\mathbf q,t_{\mathrm{sep}},W_\tau)\).

%=============================================================================
\section{Implemented Upgrade Path}
%=============================================================================

\begin{enumerate}[label=\arabic*.]
  \item \textbf{Nonzero-\(q\) validation.}
  The response schema stores
  \[
    \mathbf p_i = \mathbf p_f - \mathbf q,
  \]
  for every record.  The phase test covers both ordinary and Breit-frame
  examples, including
  \[
    \mathbf p_f=(0,0,1),\quad
    \mathbf q=(0,0,1),\quad
    \mathbf p_i=(0,0,0).
  \]
  and
  \[
    \mathbf p_f=(0,0,1),\quad
    \mathbf q=(0,0,2),\quad
    \mathbf p_i=(0,0,-1).
  \]

  \item \textbf{Multi-\(t_{\mathrm{sep}}\), \(C_2\), and summed-ratio output.}
  The application accepts a list of source-sink separations and saves
  \[
    \Ctwo(t_{\mathrm{sep}}),\quad
    \Cresp(t_{\mathrm{sep}}),\quad
    \Rsum(t_{\mathrm{sep}}).
  \]
  The saved ratios are
  \begin{equation}
    \Rsum^{\mathrm{resp}}(t_{\mathrm{sep}})
      =
      \frac{\Cresp(t_{\mathrm{sep}})}
           {\Ctwo(t_{\mathrm{sep}})},
    \qquad
    \Rsum^{\mathrm{explicit}}(t_{\mathrm{sep}})
      =
      \frac{\sum_{\tau\in W_\tau}\Cthree(t_{\mathrm{sep}},\tau)}
           {\Ctwo(t_{\mathrm{sep}})}.
  \end{equation}

  \item \textbf{Restricted tau windows.}
  The supported choices are all time slices, source-to-sink,
  open source-to-sink, a symmetric restricted window, and an explicit
  integer range:
  \[
    W_\tau=\{\tau_{\min},\ldots,t_{\mathrm{sep}}-\tau_{\min}\}.
  \]
  For restricted windows, the inserted source is summed before inversion.

  \item \textbf{Multi-current gamma response suite.}
  The application accepts multiple gamma labels from the standard pion gamma
  list.  The primary EMFF current is still \(V_4\), labeled \(\texttt{T}\),
  but scalar, pseudoscalar, vector, axial, and tensor bilinears can be used as
  current-response diagnostics without changing the first-order contraction.

  \item \textbf{HDF5 schema upgrade.}
  Version-3 first-order output stores file-level kinematics and one
  \(\texttt{results/record\_XXXX}\) group per diagnostic.  Each record includes
  \(\mathbf p_f\), \(\mathbf q\), inferred \(\mathbf p_i\), \(t_{\mathrm{sep}}\),
  current gamma, source/sink gamma labels, the tau-window definition, \(C_2\),
  explicit summed \(C_3\), response \(C_2\)-like data, ratios, and the
  response-explicit difference.  It also stores
  \texttt{source\_position}, \texttt{source\_time},
  \texttt{tau\_relative\_list}, \texttt{tau\_absolute\_list}, and
  \texttt{time\_axis=source\_relative}.  The explicit \(C_3\), response
  correlator, and \(C_2\) arrays are rolled by \(-t_0\) before analysis and
  saving.  The nested current-current diagnostic uses the same rule in its
  version-2 schema for both insertion windows.
  For analysis convenience, the same per-record scalars are also collected in
  \(\texttt{summary/}\) datasets:
  \[
    \texttt{relative\_difference},\quad
    \texttt{response\_R\_sum},\quad
    \texttt{explicit\_R\_sum},
  \]
  together with \(\mathbf p_f\), \(\mathbf q\), \(\mathbf p_i\),
  \(t_{\mathrm{sep}}\), gamma labels, and window labels.  Analysis scripts can
  therefore scan the diagnostics without opening every
  \(\texttt{results/record\_XXXX}\) group.

  \item \textbf{Optional GPU regression tests.}
  The Perlmutter application remains a small S8T32 smoke workflow.  The test
  suite contains CPU-only schema, phase, and algebra checks, while the GPU
  smoke run can be launched when a QUDA-enabled node is available.
\end{enumerate}

\paragraph{Explicit non-goal.}
This implementation does not save a per-\(\tau\) response-propagator cache.
If arbitrary post-hoc tau windows become necessary, that should be considered
as a separate design discussion rather than silently added to this diagnostic.

%=============================================================================
\section{Summary}
%=============================================================================

The present pion current background-response code is a linearized
Feynman-Hellmann diagnostic:
\begin{equation}
  \Sresp=\Dslash^{-1}\Gamma\Phi_{\mathbf q}S.
\end{equation}
It does not compute finite-\(\lambda\) energies.  Instead, it computes
\(\partial C_\lambda/\partial\lambda|_{\lambda=0}\), which equals a summed
three-point function up to the convention-dependent sign.  The energy-shift
derivative is then encoded in the slope of the response ratio
\(\Cresp/\Ctwo\) as a function of \(t_{\mathrm{sep}}\).

\end{document}
